\section{Introducció i Objectius}

La indústria farmacèutica opera sota requisits estrictes de qualitat i traçabilitat. Quan es detecta una desviació de procés o resultat fora d'especificació, cal activar una investigació de causes arrel (RCA) per identificar l'origen del problema i definir accions correctives. Malgrat la maduresa normativa del sector, aquestes investigacions continuen sent manuals i dependents de l'experiència, generant variabilitat entre investigadors, escalabilitat limitada i dificultat per garantir consistència temporal.

En aquest treball, el problema s'analitza en el context del fraccionament de plasma humà fins a l'obtenció de la Fracció II+III. Aquest procés reuneix les condicions que fan un cas d'ús especialment idoni per a ML: alta dimensionalitat (més de 150 variables de procés), variabilitat inherent de la matèria primera biològica, interaccions no lineals entre paràmetres fisicoquímics i un volum de dades històriques suficient per entrenar models robustos. A més, la seva estructura multietapa i la sensibilitat del \emph{yield} a petites desviacions operatives el converteixen en un banc de proves representatiu de processos farmacèutics complexos, de manera que la metodologia desenvolupada és transferible a altres línies de producció amb característiques similars (biotecnologia, API químics, productes estèrils). Es tracta, per tant, d'un procés farmacèutic multietapa especialment sensible a variacions fisicoquímiques i operatives. Aquesta complexitat fa que petites desviacions de procés puguin traduir-se en impactes rellevants sobre el \emph{yield} i els atributs crítics de qualitat, reforçant l'interès d'una aproximació RCA assistida per dades \cite{cohn1946,burnouf2007_modern}.

L'aparició de tècniques de Machine Learning i mètodes d'explicabilitat (XAI) permet una RCA més sistemàtica i basada en dades \cite{shap_arxiv1705,lime_arxiv1602,doshi2017rigorous}. Aquest treball proposa una metodologia d'\textbf{industrialització de la RCA basada en ML interpretable} per augmentar el criteri expert amb evidència quantitativa prioritzada, reduint l'espai d'hipòtesis i millorant la velocitat d'investigació. La proposta integra preparació de dades, modelització amb Explainable Boosting Machines (EBM), explicabilitat nativa i validació orientada a entorns regulats GxP.

\textbf{Objectius:} (1) Definir un marc RCA data-driven alineat amb requisits GxP; (2) Construir un pipeline de dades reproduïble; (3) Desenvolupar models ML interpretables per predir desviacions; (4) Aplicar tècniques d'explicabilitat per prioritzar hipòtesis de causa arrel; (5) Proposar una arquitectura d'industrialització i validació per a integració progressiva en operació.
